
\begin{exercise}[subtitle={Properties of the cumulant generating function}]
    \label{ex: properties cumulant generating function}
    Let \(X\) be a random variable. Then the cumulant generating function
    \(\CGF_X(t) = \log \MGF_X(t)\) has the following properties:
    \begin{enumerate}[label=(\roman*), noitemsep]
        \item\label{it: cgf 0=0} \(\CGF_X(0) = 0\).
        \item\label{it: cgf linear} \(\CGF_{aX + b}(t) = \scp{b, t} + \CGF_X(\scp{a, t})\) for any \(a \in \real\), \(b,t\in \hilbert\).
        \item\label{it: cumulant independent sums} Let \(Y\) be \underline{independent} from \(X\), then
        \[
            \CGF_{X+Y}(t) = \CGF_X(t) + \CGF_Y(t)
        \]
    \end{enumerate}
    If the moment generating function is finite in a neighborhood of zero,
    then \(\CGF_X\) is infinitely often differentiable and
    \begin{enumerate}[label=(\roman*), resume, noitemsep]
        \item \(\CGF_X'(0) = \E{X}\),
        \item \(\CGF_X''(0) = \var{X}\)
    \end{enumerate}
\end{exercise}

\begin{solution}
    For \ref{it: cgf 0=0}-\ref{it: cumulant independent sums} follow from the
    corresponding properties of the MGF (Exercise \ref{ex: properties moment
    generating functions}) by application of the logarithm.  For example,
    \(\CGF_X(0)=0\) follows directly from \(\MGF_X(0)=1\). For convexity
    let \(t, s\in \hilbert\) and \(\lambda \in [0,1]\). 

    The derivatives are straightforward calculations using the derivatives of the MGF
    and the chain rule due to \(\CGF_X(t) = \log \MGF_X(t)\). That is
    \begin{alignat*}{3}
        \CGF_X'(t)
        &= \frac{\MGF_X'(t)}{\MGF_X(t)}
        \qquad &\implies &\CGF_X'(0) = \E{X},
        \\
        \CGF_X''(t)
        &= \frac{\MGF_X''(t)}{\MGF_X(t)} - \frac{(\MGF_X'(t))^2}{(\MGF_X(t))^2}
        \qquad &\implies &\CGF_X''(0)
        \begin{aligned}[t]
        &= \E{X^2} - (\E{X})^2
        \\
        &= \var{X}.
        \end{aligned}
    \end{alignat*}
    This proves the claims.
\end{solution}
